\paragraph{The terrain model's inputs and fitted mechanism.}
Let $Z_t$ be a buffer of up to three visible or generated positions and
times. Its first-to-last secant gives $v_t$, its normalized direction
gives $h_t$, and $\eta_t$ indicates whether that direction is valid.
The history design is
$r(Z_t)=(1,v_{t,e},v_{t,n},\|v_t\|,h_{t,e},h_{t,n},\eta_t)^\top$.
The terrain drift and position dynamics are
\begin{equation}\label{eq:main-terrain-model}
 b_g(X_t,Z_t)=B_0^\top r(Z_t)+\theta_g^\top\widetilde\phi_g(X_t,Z_t),
 \qquad dX_t=b_g(X_t,Z_t)\,dt+L_g\,dW_t .
\end{equation}
$\widetilde\phi_g$ appends an intercept to the selected map/history
feature vector, which contains encoded numerical values and validity
masks. The road and river inputs describe nearest-feature distances and
directions; land cover uses grouped indicators; surface inputs describe
slope/orientation; registered interactions combine selected map and
history quantities. These are fixed engineering representations. Their
coefficients are learned from displacement rates, not constrained by a
law of road attraction or river avoidance.

The shared $B_0$ is fitted first by ridge regression with strength
$10^{-6}$. Each of the 404 outer-training windows contributes one actual
first-future transition $(\Delta x_i,\Delta t_i)$; the elapsed times need
not equal the method interval $\tau$. Then configuration $g$ fits the
remaining rate $y_i=\Delta x_i/\Delta t_i-B_0^\top r(Z_i)$ by
\begin{equation}\label{eq:main-terrain-fit}
 \widehat\theta_g=\underset{\theta}{\operatorname{argmin}}
 \left\{\frac1{2n}\|D_g\theta-Y\|_F^2+
             \frac\lambda2\|\theta\|_F^2\right\},
 \quad n=404,\quad\lambda=10^{-4} .
\end{equation}
$D_g$ stacks the augmented feature rows. This two-stage fit is linear in
its parameters; map lookup and history feedback nevertheless make its
position dynamics state dependent. Each configuration estimates its own
diffusion from fitted displacement residuals,
$Q_g=L_gL_g^\top=n^{-1}\sum_i
(\Delta x_i-\widehat b_{g,i}\Delta t_i)
(\Delta x_i-\widehat b_{g,i}\Delta t_i)^\top/\Delta t_i$.
The diffusion estimate differs from the method branch's centred $\tau R$.
Euler--Maruyama advances each path with drift $b_g\Delta$ and noise
$\sqrt\Delta L_gz$; every five saved seconds the history buffer appends
its generated position and recomputes the secant. Maps are queried at
the predicted position. Future observed positions never enter this buffer.

The base configuration retains the shared history drift and its fitted
history-only residual conditioner; it queries no map layer. All-terrain adds all registered
groups and interactions; LOO removes one group from that full configuration,
and LIO adds one group to the history baseline. Thus RQ2 compares these
forecasting configurations on matched targets. The shared baseline and
separate residual fits make the intended increment explicit, while
refitting and feature penalties limit attribution to an individual factor.
